\documentclass{article}
\usepackage[utf8]{inputenc}
\usepackage[spanish]{babel}
\usepackage{listings}
\usepackage{xcolor}
\usepackage{geometry}
\usepackage{amsmath}
\geometry{a4paper, margin=1in}
\definecolor{codegreen}{rgb}{0,0.6,0}
\definecolor{codegray}{rgb}{0.5,0.5,0.5}
\definecolor{codepurple}{rgb}{0.58,0,0.82}
\definecolor{backcolour}{rgb}{0.95,0.95,0.92}
\lstset{backgroundcolor=\color{backcolour},commentstyle=\color{codegreen},keywordstyle=\color{magenta},numberstyle=\tiny\color{codegray},stringstyle=\color{codepurple},basicstyle=\ttfamily\small,breakatwhitespace=false,breaklines=true,captionpos=b,keepspaces=true,numbers=left,numbersep=5pt,showspaces=false,showstringspaces=false,showtabs=false,tabsize=2,language=C++}
\newcommand{\quees}[1]{\section*{🤔 ¿QUÉ ES?}\hrulefill \\ #1}
\newcommand{\dichodeotraforma}[1]{\section*{💬 DICHO DE OTRA FORMA}\hrulefill \\ #1}
\newcommand{\diagrama}[1]{\section*{🎨 DIAGRAMA}\hrulefill \\ #1}
\newcommand{\comofunciona}[1]{\section*{⚙️ ¿CÓMO FUNCIONA?}\hrulefill \\ #1}
\newcommand{\complejidad}[1]{\section*{⏱️ COMPLEJIDAD (Big-O)}\hrulefill \\ #1}
\newcommand{\entrevistas}[1]{\section*{🎯 EN ENTREVISTAS}\hrulefill \\ #1}
\newcommand{\resumen}[1]{\section*{📋 RESUMEN RÁPIDO}\hrulefill \\ #1}
\begin{document}
\title{Sliding Window (Ventana Deslizante)}
\date{}
\maketitle
\quees{Una ventana de tren 🚂 donde ves paisajes mientras el tren avanza. 
La ventana tiene un tamaño fijo (por ejemplo, 3 metros de ancho).

Al avanzar, entra nuevo paisaje por un lado y sale el paisaje viejo
por el otro. No tienes que redibujar todo el paisaje cada vez, ¡solo
actualizas los bordes!

En programación, este patrón sirve para procesar sub-arreglos o 
sub-cadenas sin tener que recalcular todo desde cero en cada paso, 
lo cual ahorra muchísimo tiempo.}
\dichodeotraforma{El patrón Sliding Window convierte un problema de loops anidados O(n*k) 
en un problema de un solo loop O(n). 

Se usa cuando te piden algo sobre elementos CONTIGUOS de un array o 
string. Existen dos tipos:
1. Fixed Size: La ventana siempre mide K elementos.
2. Variable Size: La ventana crece o se encoge según una condición
   (usualmente usando dos punteros).}
\diagrama{\begin{verbatim}
Array: [ 1, 3, 2, 6, -1, 4 ]
Meta: Suma de cada ventana de tamaño 3.

Paso 1: [ 1, 3, 2 ] , 6, -1, 4  -> Suma = 6
Paso 2: 1, [ 3, 2, 6 ] , -1, 4  -> Suma = 11 (6 - 1 + 6)
Paso 3: 1, 3, [ 2, 6, -1 ] , 4  -> Suma = 7  (11 - 3 + (-1))

¡Solo restas el que sale y sumas el que entra! ⚡
\end{verbatim}}
\comofunciona{Fixed Size Window:
1. Calculas el resultado de la primera ventana (índices 0 a K-1).
2. Empiezas un loop desde el índice K hasta el final.
3. En cada paso:
   - Sumas el nuevo elemento (`arr[i]`).
   - Restas el elemento que sale (`arr[i - K]`).
   - Actualizas el resultado máximo/mínimo.

Variable Size Window:
1. Usas dos punteros: `start` y `end` (ambos en 0).
2. Mueves `end` para expandir la ventana y cumplir una condición.
3. Si la condición se rompe, mueves `start` para encoger la ventana.}
\section*{📝 PSEUDOCÓDIGO}
\hrulefill
\begin{verbatim}
función sumaMaxima(arr, k):
    sumaActual = suma de los primeros k elementos
    maximaSuma = sumaActual
    PARA i desde k hasta arr.tamaño - 1:
        sumaActual = sumaActual + arr[i] - arr[i - k]
        maximaSuma = max(maximaSuma, sumaActual)
    return maximaSuma
\end{verbatim}
\section*{💻 CÓDIGO C++}
\hrulefill
\begin{lstlisting}[language=C++]
#include <iostream>
#include <vector>
#include <algorithm>
using namespace std;

// Máxima suma de un subarray de tamaño K -> O(n)
int maxSubarraySum(vector<int>& arr, int k) {
    if (arr.size() < k) return -1;

    int window_sum = 0;
    for (int i = 0; i < k; i++) window_sum += arr[i];

    int max_sum = window_sum;

    for (int i = k; i < arr.size(); i++) {
        // Deslizar la ventana: sumar el nuevo, restar el viejo
        window_sum += arr[i] - arr[i - k];
        max_sum = max(max_sum, window_sum);
    }

    return max_sum;
}

int main() {
    vector<int> data = {2, 1, 5, 1, 3, 2};
    int k = 3;
    cout << "Máxima suma (k=3): " << maxSubarraySum(data, k) << endl; // 9 (5+1+3)
    return 0;
}
\end{lstlisting}
\complejidad{Caso            │ Tiempo        │ Espacio
  ────────────────┼───────────────┼────────────────────────────────
  Sliding Window  │ O(n)          │ O(1)
  Fuerza Bruta    │ O(n * k)      │ O(1)

* La diferencia es abismal cuando `n` y `k` son grandes.}
\entrevistas{}
\resumen{• Elementos contiguos.
• Evita recalcular: Actualiza los bordes de la ventana.
• Reduce O(n*k) a O(n).
• Dos tipos: Fija y Variable.
• Ideal para sumas de rango, promedios móviles y búsqueda de sub-strings.
════════════════════════════════════════════════════════════════}
\end{document}